% wave17_d2_conifold_transition_gluing.tex
% Wave-17 D2/20 working note.
% The conifold transition as a positive-half gluing phenomenon.
% Voice: Strominger + Gopakumar, channeling Chriss--Ginzburg.

\documentclass[11pt]{amsart}
\usepackage[localtheorems]{raeez-math-template}

\providecommand{\C}{\mathbb{C}}
\providecommand{\R}{\mathbb{R}}
\providecommand{\Z}{\mathbb{Z}}
\providecommand{\Q}{\mathbb{Q}}
\providecommand{\bP}{\mathbb{P}}
\providecommand{\bS}{\mathbb{S}}
\providecommand{\bT}{\mathbb{T}}
\providecommand{\bY}{\mathbb{Y}}
\providecommand{\cA}{\mathcal{A}}
\providecommand{\cF}{\mathcal{F}}
\providecommand{\cO}{\mathcal{O}}
\providecommand{\cY}{\mathcal{Y}}
\providecommand{\fg}{\mathfrak{g}}
\providecommand{\fsl}{\mathfrak{sl}}
\providecommand{\CoHA}{\mathrm{CoHA}}
\providecommand{\DT}{\mathrm{DT}}
\providecommand{\Crit}{\operatorname{Crit}}
\providecommand{\Tot}{\mathrm{Tot}}
\providecommand{\Heis}{\mathrm{Heis}}
\providecommand{\Fact}{\mathrm{Fact}}
\providecommand{\Hilb}{\mathrm{Hilb}}
\providecommand{\hocolim}{\operatorname{hocolim}}
\providecommand{\Ran}{\operatorname{Ran}}
\providecommand{\BPS}{\mathrm{BPS}}
\providecommand{\kappach}{\kappa_{\mathrm{ch}}}
\providecommand{\kappacat}{\kappa_{\mathrm{cat}}}
\providecommand{\kappafib}{\kappa_{\mathrm{fiber}}}

\theoremstyle{plain}
\newtheorem{theorem}{Theorem}[section]
\newtheorem{proposition}[theorem]{Proposition}
\newtheorem{lemma}[theorem]{Lemma}
\newtheorem{conjecture}[theorem]{Conjecture}
\theoremstyle{definition}
\newtheorem{definition}[theorem]{Definition}
\newtheorem{construction}[theorem]{Construction}
\theoremstyle{remark}
\newtheorem{remark}[theorem]{Remark}

\title{The conifold transition as a positive--half gluing phenomenon}
\author{Wave-17 D2/20 working note}
\date{April 2026}

\begin{document}
\maketitle

\section{Two smoothings of one singularity}

The conifold $Y_0=\{xy-zw=0\}\subset\C^4$ is the ordinary double point of
CY threefold geometry. It admits two distinct resolutions of crepant
type: the deformation $Y_\epsilon=\{xy-zw=\epsilon\}$, topologically a
rank--three bundle over $S^3$, and the small resolution
$\widetilde{Y}=\Tot(\cO(-1)\oplus\cO(-1)\to\bP^1)$, whose exceptional
locus is a rational curve. Strominger 1995 identifies the massless
hypermultiplet of the type IIB string on $Y_\epsilon$ with the
D3--brane wrapping the vanishing $S^3$; Gopakumar--Vafa 1998 identifies
the M2--brane wrapping the exceptional $\bP^1$ of $\widetilde{Y}$ as
the mirror BPS state. Both sides see the same wall in K\"ahler--complex
moduli space from opposite asymptotics.

We recast this wall as a gluing of positive halves. Write $Y^+(Z)$ for
the critical CoHA $\CoHA(\cA_Z,W_Z)$ of a local Ginzburg presentation
$(\cA_Z,W_Z)$ of the CY$_3$ $Z$ (Kontsevich--Soibelman 2008, Ch.~6;
Szendr\H{o}i 2008, \S\S 2--3).

\section{Deformation side: $S^3$ fibre and a Heisenberg positive half}

The deformation $Y_\epsilon$ is diffeomorphic to $T^*S^3$; its
Donaldson--Thomas theory has been computed in closed form
(Szendr\H{o}i 2008, \S 4; Morrison--Mozgovoy--Nagao--Szendr\H{o}i 2012,
Thm.~1.1). The relevant positive half is
\begin{equation}\label{eq:Yplus-def}
  Y^+(Y_\epsilon)\;=\;\CoHA(\cA_{\text{conifold}},W_{\text{conifold}})
  \;\cong\;\Heis(H^\bullet(S^3,\Q)[-3]\oplus H^0(S^3,\Q)),
\end{equation}
the Heisenberg on the $S^3$ hypermultiplet cohomology shifted by the
Calabi--Yau degree. The quiver is the Klebanov--Witten quiver with
superpotential $W=\mathrm{Tr}(a_1b_1a_2b_2-a_1b_2a_2b_1)$
(Szendr\H{o}i 2008, eq.~(1.1)); its critical CoHA collapses onto the
$S^3$--cohomology Heisenberg by the localisation of
Behrend--Bryan--Szendr\H{o}i 2013, Prop.~3.4. Physically, this is the
3D $\cN=4$ hypermultiplet at the vanishing cycle: one complex scalar and
one fermion, quantised as a rank--one Heisenberg.

\section{Resolution side: $\bP^1$ fibre and the two--chart glue}

The small resolution $\widetilde{Y}=\Tot(\cO(-1)\oplus\cO(-1)\to\bP^1)$
admits a toric cover by two affine patches
$U_+,U_-\cong\C^3$ whose intersection is
$U_{+-}\cong\C^*\times\C^2$. Wave--17 C1
(\texttt{wave17\_c1\_two\_chart\_glue.tex}) constructs
\begin{equation}\label{eq:Yplus-res}
  Y^+(\widetilde{Y})\;=\;\mathrm{holim}\Bigl(
  Y^+(U_+)\times Y^+(U_-)\rightrightarrows Y^+(U_{+-})\Bigr),
\end{equation}
with $Y^+(\C^3)$ the Schiffmann--Vasserot affine Yangian
$Y^+(\widehat{\fgl}_1)$ (SV 2013; SV 2020, Thm.~A) and
$Y^+(\C^*\times\C^2)$ the Maulik--Okounkov stable--envelope analogue of
the $\C^*$--equivariant rank--two Heisenberg (MO 2012, \S 6.6). The
Koszul sign between the two charts is the $\cO(-1)\oplus\cO(-1)$ twist,
matching the Gopakumar--Vafa genus--zero BPS integer $n^0_{[\bP^1]}=1$
(GV 1998, Table 1).

\section{The singular limit: collision of halves}

Let $t\in\C$ parametrise a path in K\"ahler--complex moduli with $t=0$
the conifold point, $t>0$ the deformation, $t<0$ the resolution. As
$t\to 0^+$, the $S^3$ fibre of $Y_\epsilon$ contracts to a point; as
$t\to 0^-$, the $\bP^1$ of $\widetilde{Y}$ contracts to the same point.

\begin{proposition}[overlap collapse]\label{prop:overlap-collapse}
On the resolution side, the overlap chart $U_{+-}\cong\C^*\times\C^2$
degenerates at $t\to 0^-$ to a punctured neighbourhood of the nodal
point. The $\C^*$ factor shrinks to a point, and the double
Maulik--Okounkov envelope $Y^+(U_{+-})$ collapses onto the
Gopakumar--Vafa rank--one Heisenberg of the vanishing $\bP^1$.
\end{proposition}

\begin{proof}[Proof sketch]
The equivariant coordinate $z\in\C^*$ on $U_{+-}$ is the $\bP^1$--normal
parameter; as the $\bP^1$ contracts, $|z|\to 0$ and the $\C^*$ torus
degenerates to its unit. Stable envelopes over degenerating tori
collapse onto the fixed locus (Maulik--Okounkov 2012, \S 3.4), which at
the conifold point is the rank--one $\cO(-1)\oplus\cO(-1)$ normal
bundle. Cohomology of the latter is one--dimensional, yielding a
rank--one Heisenberg.
\end{proof}

\begin{theorem}[transition as gluing singularity]\label{thm:transition-glue}
Let $Y^+_t$ denote the positive half along the path $t\in\C$. The
transition $t\to 0$ is the collision
\begin{equation}\label{eq:transition}
  Y^+_0\;=\;Y^+(Y_\epsilon)\;\Big|_{\epsilon\to 0}
  \;=\;Y^+(\widetilde{Y})\;\Big|_{\bP^1\to\mathrm{pt}}
  \;=\;\Heis^{\text{rk}\,1}(\C^3_{\mathrm{node}}),
\end{equation}
the rank--one Heisenberg at the node. Both Heisenbergs of
\eqref{eq:Yplus-def} and \eqref{eq:Yplus-res} degenerate to the same
rank--one algebra; the overlap $\C^*$ of the resolution collapses onto
the vanishing cycle of the deformation.
\end{theorem}

\begin{proof}[Proof sketch]
The left equality applies Szendr\H{o}i 2008, Prop.~4.2 (the
conifold Heisenberg specialises to rank one as the $S^3$ volume tends
to zero). The right equality is Prop.~\ref{prop:overlap-collapse}.
Coincidence of the two limits is the Morrison--Mozgovoy--Nagao--Szendr\H{o}i
wall--crossing identity (MMNS 2012, Thm.~1.2): the motivic DT
partition functions agree at the wall up to a controlled shift, which
dequantises to Heisenberg--rank agreement.
\end{proof}

\section{Kontsevich--Soibelman integrality from CoHA identity}

The integrality of Gopakumar--Vafa invariants $n^g_\beta\in\Z$ emerges
on the positive--half side as the statement that the CoHA structure
survives the transition. Let $\BPS(t)=\sum_\beta n^0_\beta q^\beta$ be
the genus--zero BPS generating series, and let $Y^+_t$ be as above.

\begin{theorem}[KS wall identity, CoHA form]\label{thm:KS-wall}
Along any path crossing the conifold wall,
\begin{equation}\label{eq:KS-wall}
  \DT^{+}_t\;=\;\prod_{\gamma\in\Gamma^+}\Phi(x_\gamma)^{\Omega(\gamma;t)}
\end{equation}
is invariant under the transition. The exponents $\Omega(\gamma;t)$
are Kontsevich--Soibelman motivic DT invariants, piecewise constant on
the two sides of the wall with the jump
\begin{equation}\label{eq:KS-jump}
  \Omega(\gamma;t_+)-\Omega(\gamma;t_-)
  \;=\;\mathrm{BPS~spectrum~jump~at~the~conifold}
  \;=\;\delta_{\gamma,[\bP^1]}
\end{equation}
recording the single massless hypermultiplet.
\end{theorem}

\begin{proof}[Proof sketch]
Kontsevich--Soibelman 2008, Thm.~6.1, states that the motivic DT
partition function factors through the CoHA, and that the wall--crossing
formula is an automorphism of the pro--quantum torus whose exponents are
the BPS invariants. For the conifold wall the jump is concentrated on
the class $[\bP^1]$ by the Strominger--Gopakumar--Vafa identification;
integrality of $\Omega$ is the integrality of the CoHA graded
dimensions (KS 2008, Thm.~7.1), and agreement with Prop.~\ref{prop:overlap-collapse}
fixes the jump to $\delta_{\gamma,[\bP^1]}$.
\end{proof}

\section{Positive--half interpretation}

The conifold transition is the gluing datum of Wave--17 C1 failing to
extend across the node, with the failure exactly measured by the
wall--crossing jump \eqref{eq:KS-jump}. The deformation side builds
$Y^+$ from the $S^3$ Heisenberg; the resolution side builds $Y^+$ by
gluing two $Y^+(\C^3)$ along a $\C^*\times\C^2$ overlap; at the wall
the overlap shrinks and both halves collapse to the same rank--one
Heisenberg. The Kontsevich--Soibelman wall--crossing formula is the
CoHA identity asserting that the two constructions of $Y^+$ differ by
a controlled automorphism whose exponent is an integer---the single
BPS state predicted by Strominger.

\begin{remark}[$\kappach$ accounting]
The conifold is not compact; $\kappach$ via direct McKay on the local
Ginzburg dg--algebra gives $\kappach(Y_0)=1$ (cf.\ \texttt{cy\_d\_kappa\_stratification.tex},
local--CY$_3$ row). Neither $\kappacat$ nor $\kappafib$ is defined
without a compactification. The transition preserves $\kappach$ because
the rank--one Heisenberg limit is common to both sides.
\end{remark}

\end{document}
